\documentclass[../../main.tex]{subfiles}
\graphicspath{{../../image/}{../../figure/}}
\begin{document}

\section{总变差的界}\label{section:ch14-总变差的界}

如\cref{section:ch13-glimm泛函}，$TV_{(-\infty,\infty)}U_\delta(\cdot,t)$ 将通过
\begin{align}
L(t)=\sum|\gamma|,
\label{align:ch14-L}
\end{align}
即穿过 $t$-时间线的所有跳跃间断的强度之和来度量。显然，$L(\cdot)$ 在相邻前沿碰撞之间的时间区间上保持常数，仅在波相互作用点改变。为估计这些变化，须考察各类碰撞。

设振幅为 $\alpha$ 的 $j$-前沿与振幅为 $\beta$ 的 $i$-前沿碰撞。当 $|\alpha||\beta|\geqslant\sigma$ 时，所得跳跃间断经近似 Riemann 解算器分解为完整波扇 $\varepsilon=(\varepsilon_1,\cdots,\varepsilon_n)$，于是由\cref{theorem:ch09-波扇相互作用}，
\begin{align}
|\varepsilon_j-\alpha|+|\varepsilon_i-\beta|+\sum_{k\neq i,j}|\varepsilon_k|=O(1)|\alpha||\beta|,
\label{align:ch14-interaction-ij}
\end{align}
当 $i<j$；或
\begin{align}
|\varepsilon_i-\alpha-\beta|+\sum_{k\neq i}|\varepsilon_k|=O(1)|\alpha||\beta|,
\label{align:ch14-interaction-ii}
\end{align}
当 $i=j$。另一方面，当 $|\alpha||\beta|<\sigma$ 时，所得跳跃间断经简化 Riemann 解算器求解（如\cref{figure:ch14-不同族前沿碰撞}或\cref{figure:ch14-同族前沿碰撞}所示），碰撞前沿的振幅守恒。生成出射伪激波的强度可由状态空间中的波图容易地估计：
\begin{align}
|U_R-U_Q|=O(1)|\alpha||\beta|.
\label{align:ch14-pseudoshock-gen}
\end{align}

接着考虑\cref{figure:ch14-伪激波与前沿碰撞}所示情形，其中伪激波与振幅为 $\beta$ 的 $i$-前沿碰撞。由于 $i$-前沿的振幅跨碰撞守恒，对状态空间波图（\cref{figure:ch14-伪激波与前沿碰撞}）的分析给出：出射伪激波的强度与入射伪激波的强度有关系
\begin{align}
|U_R-U_Q|=|U_M-U_L|+O(1)|\beta||U_M-U_L|.
\label{align:ch14-pseudoshock-hit}
\end{align}

令 $I$ 表示发生碰撞的 $t\in(0,T)$ 集合。记 $\Delta$ 为 $t\in I$ 时从 $t-$ 到 $t+$ 的``跳跃''算子。根据上述波相互作用分析，我们推断
\begin{align}
\Delta L(t)\leqslant\kappa|\alpha||\beta|,\qquad t\in I,
\label{align:ch14-deltaL}
\end{align}
其中 $|\alpha|$ 与 $|\beta|$ 是在 $t$ 处碰撞的波的强度。

我们把 $TV_{(-\infty,\infty)}U_\delta(\cdot,t)$ 控制住的策略是：证明\eqref{align:ch14-deltaL}所允许的 $L(\cdot)$ 任一增量，都被同时减少的潜在波相互作用量所抵消。

一个 $j$-波与一个 $i$-波，若前者穿过 $t$-时间线位于后者之左，则当 $i<j$ 或 $i=j$ 且至少一个波是压缩激波时，称它们是\textbf{接近的}。

在 $t\in(0,T)\setminus I$ 处的波相互作用势由
\begin{align}
Q(t)=\sum|\zeta||\xi|,\qquad t\in(0,T)\setminus I,
\label{align:ch14-Q}
\end{align}
度量，其中求和遍历所有接近波对（其强度记为 $|\zeta|$ 与 $|\xi|$），它们穿过 $t$-时间线。特别地，
\begin{align}
Q(t)\leqslant\tfrac12 L(t)^2,\qquad t\in(0,T)\setminus I.
\label{align:ch14-Q-bound}
\end{align}
显然 $Q(\cdot)$ 在相邻碰撞之间保持常数。另一方面，在任一 $t\in I$，强度 $|\alpha|$ 与 $|\beta|$ 的波碰撞，波相互作用分析蕴含
\begin{align}
\Delta Q(t)\leqslant-|\alpha||\beta|+\kappa|\alpha||\beta|\,L(t-),\qquad t\in I.
\label{align:ch14-deltaQ}
\end{align}

类比 Glimm 泛函\eqref{align:ch13-G-def}，置
\begin{align}
G(t)=L(t)+2\kappa Q(t),\qquad t\in(0,T)\setminus I.
\label{align:ch14-G}
\end{align}

结合\eqref{align:ch14-G}、\eqref{align:ch14-deltaL}与\eqref{align:ch14-deltaQ}得
\begin{align}
\Delta G(t)\leqslant\kappa[2\kappa G(t-)-1]|\alpha||\beta|,\qquad t\in(0,T)\setminus I.
\label{align:ch14-deltaG}
\end{align}
设初值总变差充分小，使 $4\kappa L(0+)\leqslant1$。于是由\eqref{align:ch14-G}与\eqref{align:ch14-Q-bound}，$G(0+)\leqslant2L(0+)\leqslant(2\kappa)^{-1}$。连同\eqref{align:ch14-deltaG}与简单归纳论证给出 $\Delta G(t)\leqslant0$，$t\in I$，即 $G(\cdot)$ 非增。因而
\begin{align}
L(t)\leqslant G(t)\leqslant G(0+)\leqslant2L(0+),\qquad t\in(0,T)\setminus I,
\label{align:ch14-L-bound}
\end{align}
这建立了所需估计\eqref{align:ch14-delta-tv}。

接着估计波相互作用总量。由于 $\kappa L(t-)\leqslant\frac12$，\eqref{align:ch14-deltaQ}给出
\begin{align}
\Delta Q(t)\leqslant-\tfrac12|\alpha||\beta|,\qquad t\in I.
\label{align:ch14-deltaQ-half}
\end{align}
对所有 $t\in I$ 求和\eqref{align:ch14-deltaQ-half}，并利用\eqref{align:ch14-Q-bound}，
\begin{align}
\sum|\alpha||\beta|\leqslant L(0+)^2,
\label{align:ch14-interaction-total}
\end{align}
其中求和遍历 $(-\infty,\infty)\times(0,T)$ 中的碰撞集合。

现考虑 $(-\infty,\infty)\times[0,T)$ 中任一 Lipschitz 曲线 $\mathscr{C}$，它关于 $U_\delta$ 非共振，定义见\cref{section:ch14-全局波型}。目的是估计 $U_\delta$ 在 $\mathscr{C}$ 上迹的总变差，它由撞上 $\mathscr{C}$ 的所有波强度之和 $L_{\mathscr{C}}=\sum|\gamma|$ 度量。

令 $J$ 表示有波撞上 $\mathscr{C}$ 的 $t\in(0,T)$ 集合。对 $t\in(0,T)\setminus(I\cup J)$，置
\begin{align}
M(t)=\sum_-|\gamma|+\sum_+|\gamma|+\sum_0|\gamma|,
\label{align:ch14-M}
\end{align}
其中 $\sum_-$（或 $\sum_+$）遍历在 $\mathscr{C}$ 正侧（或负侧）穿过 $t$-时间线的 $i$-波（$i\in\mathscr{N}_-$，相应地 $i\in\mathscr{N}_+$）；$\sum_0$ 遍历在 $\mathscr{C}$ 任一侧穿过 $t$-时间线的 $i$-波（$i\in\mathscr{N}_0$）。显然，
\begin{align}
\Delta M(t)=-|\gamma|,\qquad t\in J\setminus I,
\label{align:ch14-deltaM1}
\end{align}
\begin{align}
\Delta M(t)\leqslant\kappa|\alpha||\beta|,\qquad t\in I\setminus J,
\label{align:ch14-deltaM2}
\end{align}
\begin{align}
\Delta M(t)\leqslant-|\gamma|+\kappa|\alpha||\beta|,\qquad t\in I\cap J,
\label{align:ch14-deltaM3}
\end{align}
其中 $|\alpha|$ 与 $|\beta|$ 是在 $t\in I$ 碰撞的波强度，$|\gamma|$ 是在 $t\in J$ 撞上 $\mathscr{C}$ 的波强度。对所有 $t\in I\cup J$ 求和上述不等式，并利用\eqref{align:ch14-interaction-total}与 $4\kappa L(0+)\leqslant1$，我们得到
\begin{align}
L_{\mathscr{C}}\leqslant M(0+)+\kappa\sum|\alpha||\beta|\leqslant2L(0+).
\label{align:ch14-LC-bound}
\end{align}

波相互作用量有界的另一个重要推论是碰撞总数有限且有界（与 $T$ 无关）。确实，回忆近似 Riemann 解算器仅在两入射前沿强度之积超过 $\sigma$ 时才用于求解碰撞。由\eqref{align:ch14-interaction-total}，此类碰撞的数目以 $L(0+)^2/\sigma$ 为界。前沿仅由应用近似 Riemann 解算器求解 $U_{0\delta}$ 的跳跃间断或前沿碰撞而生成。因而前沿数目有界。任两前沿在其生存期内至多碰撞一次，故前沿间碰撞数也有界。由于所有伪激波都由前沿碰撞生成，伪激波数目同样有界。但进而前沿与伪激波间的碰撞数也必有界。总之，碰撞总数有限，其界仅依赖于 $\delta,\sigma$ 与 $U_{0\delta}$ 的跳跃点数。因而前沿追踪算法在整个上半平面上经有限步生成 $U_\delta$。特别地，估计\eqref{align:ch14-delta-tv}与\eqref{align:ch14-delta-lip}分别对 $0\leqslant t<\infty$ 与 $0\leqslant\tau<t<\infty$ 成立。

\end{document}
